\textbf{结论名称：} 函数奇偶性判定步骤（定义域对称优先）

\textbf{适用条件：}
函数给出的是解析式形式，需要快速判断奇偶性。
第一步必须检查定义域是否关于原点对称（任何非对称都直接得出非奇非偶）。

\textbf{变量说明：}
\( f(x) \) 为给定的一元实函数，定义域为 \( D \)；
原点对称指：若 \( x \in D \)，则必有 \( -x \in D \)。

\textbf{核心公式：}
定义域对称是奇偶性成立的必要条件：
\[
\boxed{\text{奇偶性成立} \;\Longrightarrow\; \forall x \in D,\; (-x) \in D}
\]
判定关系式（在对称定义域下继续验证）：
\[
\boxed{
\begin{aligned}
\text{偶函数:}&\quad f(-x) = f(x) \\
\text{奇函数:}&\quad f(-x) = -f(x)
\end{aligned}
}
\]

\textbf{使用场景：}
- 判断初等复合函数的奇偶性
- 利用奇偶性简化定积分、求值、作图
- 抽象函数性质推断时的前置检查

\textbf{简单例子：}
判断 \( f(x) = |x| + x^3 \) 的奇偶性。
解：定义域 \( \mathbb{R} \) 关于原点对称。
计算 \( f(-x) = |-x| + (-x)^3 = |x| - x^3 \)。
比较得 \( f(-x) \neq f(x) \) 且 \( f(-x) \neq -f(x) \)，故该函数为非奇非偶函数。
（注：若仅看 \( |x| \) 为偶、\( x^3 \) 为奇，复合后不一定保持单一奇偶性，需按定义验证。）